\scanpage{6}

$j<\omega$ est grand devant $\mathrm{card}\,\mathrm{Fl}\,\Theta$, pour tt $\Theta$ comme dessus. On va établir d'abord \deleted{\unsure{à}} le point suivant, plus fort que $2^\circ$ :

\[
(*)\quad\forall j<\omega,\qquad C_j\unsure{-|\to}C_\omega\text{ est }\delta_j\text{-exact.}
\]

Cela résulte aisément de a) et c) \inserted{(et grandeur de $\omega$)} / \uline{Détaillons la démonstration}

On \deleted{\unsure{puisse montrer|peut}} \inserted{en conclut} $1^\circ$), utilisant $(*)$ et b)

Pour prouver $3^\circ$, on note que \inserted{b) et} c*) [et la grandeur de $\omega$] impliquent
\[
(**)\quad\delta_\omega=\deleted{\unsure{\text{diagr.}}}(\text{cônes exacts de types }d'\text{ et }d'')\cup\alpha_\omega(\delta).
\]
\unsure{$\cup$}

Tout revient donc à construire $K(C,\delta)$, de façon à satisfaire a) : d). On construira \inserted{$(C,\delta)\longmapsto$} $K(C,\delta)$ comme composé $(C,\delta)\longmapsto\sigma\rho(C,\delta)$ de deux constructions, dont chacune fait environ la moitié du boulot :

\inserted{Dans chacune, on doit vérifier l'analogue de c) et d) (qui sera vrai pour le composé).}

$(C,\delta)\longmapsto\rho(C,\delta)$ se fait en plusieurs étapes. On considère $\left\{\begin{aligned}\Delta^{\unsure{\pi}}_{(\varnothing,d'')}&=\Delta''\\\Delta_{(d',\varnothing)}&=\Delta'\end{aligned}\right.$

\deleted{\unsure{$\rho(C,\delta)$ :}}
\[
\begin{aligned}
\rho_1(C,\delta)&=(C_1,\delta_1),\quad\text{avec }C_1\text{ la }\Delta''\text{-enveloppe de }C,\quad\deleted{\unsure{\rho_1}}\inserted{\tau_0}:C\to C_1\\
\delta'_1&=\tau_0(\delta'),\quad\delta''_1=\tau_0(\delta'')\cup(\text{diagrammes de type }d''\text{ de }C_1)
\end{aligned}
\]
\uline{Cela assure} b$'$)\\
c) et d) vérifiés

$\rho_2(C_1,\unsure{d_1|\delta_1})=(C_2,\delta_2)$, \inserted{$\tau_1:C_1-C_2$} avec $C_2$ déduit de $C_1$ en égalisant, tout couple de flèches $z\overset{u}{\underset{v}{\rightrightarrows}}x$ dans $C_1$, de but le sommet $x$ d'un $(x\xrightarrow{p_\lambda}x_\lambda)\in\delta''_1$, satisfaisant $p_\lambda u=p_\lambda v$ pour tt $\lambda$. On \unsure{prend|pose|a} $\delta_2=\tau_1(\delta_1)$

\uline{Cela assure} $a'_1$. c) et d) sont vérifiés.

$\rho_3(C_2,\delta_2)=(C_3,\delta_3)$, \inserted{$\tau_2:C_2-C_3$} avec $C_3$ déduit de $C_2$ en rajoutant, pour $\forall$ \deleted{\unsure{$u$}} \inserted{$\Theta''=$} $(x\overset{p_\lambda}{-}x_\lambda)\in\delta''_2$ et toute \uline{augmentation} $q=(z\xrightarrow{q_\lambda}x_\lambda)_\lambda$ du système projectif $(x_\lambda)$, une flèche $u_{\Theta'',q}$ à $C_3$, satisfaisant $p_\lambda u_{\Theta'',q}=q_\lambda\quad\forall\lambda$.
\[
\delta_3=\tau_2(\delta_2)
\]
\uline{Cela assure} $a'_2$\\
c) et d) vérifiés

% archive sheet number: 5
